% Part I — Title, Introduction, and Motivation (Slides 1–20)

%% Slide 1: Title
\begin{frame}[plain]
\maketitle
\end{frame}

%% Slide 2: What This Talk Is About
\begin{frame}
\frametitle{What This Talk Is About}
\begin{center}
\Large
\textit{``By the end, you'll understand a new approach to proving code correct ---\\[6pt]
one that works \jugmark{with} Python, not against it.''}
\end{center}
\vfill
\begin{itemize}
  \item No category theory required
  \item No sheaf theory required
  \item Just Python and basic logic
\end{itemize}
\end{frame}

%% Slide 3: The Stakes — Why Correctness Matters
\begin{frame}
\frametitle{The Stakes --- Why Correctness Matters}
\begin{columns}[T]
\column{0.33\textwidth}
\begin{alertblock}{Ariane 5 (1996)}
  \begin{itemize}
    \item Integer overflow
    \item 64-bit $\to$ 16-bit cast
    \item \textcolor{jugeoRed}{\textbf{\$370M}} lost in 37\,s
  \end{itemize}
\end{alertblock}
\column{0.33\textwidth}
\begin{alertblock}{Heartbleed (2014)}
  \begin{itemize}
    \item Missing bounds check
    \item OpenSSL buffer over-read
    \item \textcolor{jugeoRed}{\textbf{500K+}} servers exposed
  \end{itemize}
\end{alertblock}
\column{0.33\textwidth}
\begin{alertblock}{Therac-25 (1985)}
  \begin{itemize}
    \item Race condition
    \item Radiation therapy machine
    \item \textcolor{jugeoRed}{\textbf{6 patients}} killed or injured
  \end{itemize}
\end{alertblock}
\end{columns}
\vfill
\begin{center}
\large Bugs cost \textbf{lives} and \textbf{money}. Can we do better?
\end{center}
\end{frame}

%% Slide 4: How Do We Know Code Is Correct?
\begin{frame}
\frametitle{How Do We Know Code Is Correct?}
\begin{center}
\begin{tikzpicture}[
  box/.style={draw, rounded corners, minimum width=2.4cm, minimum height=0.9cm,
              font=\small, text=white, thick},
  arr/.style={-{Stealth[length=3mm]}, thick, jugeoBlue}
]
  \node[box, fill=jugeoRed!70]                        (rev)  {Manual Review};
  \node[box, fill=jugeoOrange, right=0.6cm of rev]    (test) {Testing};
  \node[box, fill=jugeoOrange!80, right=0.6cm of test](type) {Type Checking};
  \node[box, fill=jugeoGreen!80, right=0.6cm of type] (fv)   {Formal Verification};

  \draw[arr] (rev)  -- (test);
  \draw[arr] (test) -- (type);
  \draw[arr] (type) -- (fv);

  \node[below=0.3cm of rev,  font=\scriptsize, text=jugeoRed!70]  {Low confidence};
  \node[below=0.3cm of fv,   font=\scriptsize, text=jugeoGreen!80]{Mathematical proof};
  \node[above=0.3cm of test, font=\scriptsize, text=jugeoBlue]    {More confidence $\longrightarrow$};
\end{tikzpicture}
\end{center}
\vfill
\begin{itemize}
  \item Each step catches \emph{more} bugs --- but demands \emph{more} effort
  \item Most Python projects stop at testing
  \item \jugmark{What would it take to go further?}
\end{itemize}
\end{frame}

%% Slide 5: What Is Formal Verification?
\begin{frame}
\frametitle{What Is Formal Verification?}
\begin{block}{Definition}
A \textbf{mathematical proof} that code meets its specification --- for \emph{all} possible inputs.
\end{block}
\vspace{0.4cm}
\begin{columns}[T]
\column{0.48\textwidth}
\begin{exampleblock}{What testing says}
  ``It didn't crash on \textbf{my} tests.''
\end{exampleblock}
\column{0.48\textwidth}
\begin{exampleblock}{What verification says}
  ``It \textbf{cannot} have this class of bug.''
\end{exampleblock}
\end{columns}
\vfill
\begin{itemize}
  \item Covers \emph{infinite} input space
  \item Eliminates entire bug categories by construction
  \item The gold standard --- but historically very expensive
\end{itemize}
\end{frame}

%% Slide 6: The Big Three — LEAN, Coq, F*
\begin{frame}
\frametitle{The Big Three --- LEAN, Coq, F*}
\begin{description}[\textbf{LEAN\,4}]
  \item[\textbf{LEAN\,4}] Interactive theorem prover from \textcolor{jugeoBlue}{Microsoft Research}.\\
    Dependent types, tactic-based proofs, growing math library.
  \item[\textbf{Coq}] Foundational proof assistant from \textcolor{jugeoBlue}{INRIA} (France).\\
    Calculus of Inductive Constructions, used in CompCert \& seL4.
  \item[\textbf{F*}] Effectful verification language from \textcolor{jugeoBlue}{MSR}.\\
    Refinement types + SMT (Z3), used in Project Everest (TLS).
\end{description}
\vfill
\begin{center}
  \small All three are \textbf{powerful}, \textbf{mature}, and \textbf{not Python}.
\end{center}
\end{frame}

%% Slide 7: How They Work — The Curry-Howard Idea
\begin{frame}
\frametitle{The Curry-Howard Idea}
\begin{center}
\Large ``Proofs are programs. Programs are proofs.''
\end{center}
\vspace{0.3cm}
\begin{columns}[T]
\column{0.45\textwidth}
\begin{block}{Logic side}
  \begin{itemize}
    \item Proposition $P$
    \item Proof of $P$
    \item $P \Rightarrow Q$
  \end{itemize}
\end{block}
\column{0.1\textwidth}
\begin{center}
\vspace{0.6cm}
\Large $\cong$
\end{center}
\column{0.45\textwidth}
\begin{block}{Programming side}
  \begin{itemize}
    \item Type \texttt{P}
    \item Value of type \texttt{P}
    \item Function \texttt{P -> Q}
  \end{itemize}
\end{block}
\end{columns}
\vfill
\begin{itemize}
  \item A proof of $P$ is a \textbf{program} that has type $P$
  \item If the program type-checks, the proof is \textbf{valid}
  \item Types \emph{enforce} correctness at compile time
\end{itemize}
\end{frame}

%% Slide 8: The Extraction Model
\begin{frame}
\frametitle{The Extraction Model}
\begin{center}
\begin{tikzpicture}[
  box/.style={draw, thick, rounded corners, minimum width=2.8cm, minimum height=1cm,
              font=\small, align=center},
  arr/.style={-{Stealth[length=3mm]}, thick}
]
  \node[box, fill=jugeoLight]                     (write) {Write program\\in proof language};
  \node[box, fill=jugeoLight, right=1.5cm of write]  (prove) {Prove it\\correct};
  \node[box, fill=jugeoGreen!20, right=1.5cm of prove]  (extract) {Extract to\\``real'' language};

  \draw[arr, jugeoBlue] (write)   -- node[above, font=\scriptsize]{encode} (prove);
  \draw[arr, jugeoGreen] (prove)  -- node[above, font=\scriptsize]{extract} (extract);

  \node[below=0.6cm of write,   font=\scriptsize, text=gray] {LEAN / Coq / F*};
  \node[below=0.6cm of extract, font=\scriptsize, text=gray] {C, OCaml, Haskell};

  \node[below=0.15cm of prove, font=\scriptsize, text=jugeoBlue] {$\checkmark$ proven};
\end{tikzpicture}
\end{center}
\vfill
\begin{itemize}
  \item You re-implement your algorithm inside the proof assistant
  \item Prove properties about that implementation
  \item Extract a certified executable --- in a language \emph{they} choose
  \item \jugmark{Your original code is never verified directly}
\end{itemize}
\end{frame}

%% Slide 9: Pain Point 1 — No Python
\begin{frame}
\frametitle{Pain Point 1 --- No Python}
\begin{center}
\begin{tabular}{l l}
  \toprule
  \textbf{Proof Assistant} & \textbf{Extraction Targets} \\
  \midrule
  LEAN 4  & C \\
  Coq     & OCaml, Haskell \\
  F*      & OCaml, F\#, C, Wasm \\
  \midrule
  \multicolumn{2}{c}{\textcolor{jugeoRed}{\textbf{Python?  Nowhere.}}} \\
  \bottomrule
\end{tabular}
\end{center}
\vfill
\begin{itemize}
  \item Python is the \textbf{\#1 language} (TIOBE, Stack Overflow, GitHub)
  \item ML, data science, web backends, scripting --- all Python
  \item Formal verification simply \textbf{ignores} the largest developer community
\end{itemize}
\end{frame}

%% Slide 10: Pain Point 2 — Rewrite Everything
\begin{frame}[fragile]
\frametitle{Pain Point 2 --- Rewrite Everything}
\begin{columns}[T]
\column{0.48\textwidth}
\begin{block}{Your Python}
\begin{lstlisting}
def merge_sort(xs):
    if len(xs) <= 1:
        return xs
    mid = len(xs) // 2
    left = merge_sort(xs[:mid])
    right = merge_sort(xs[mid:])
    return merge(left, right)
\end{lstlisting}
\end{block}
\column{0.48\textwidth}
\begin{block}{Rewritten in Coq}
\begin{lstlisting}[language=ML,
  morekeywords={Fixpoint,match,with,end,
    Defined,Program,let,in}]
Fixpoint merge_sort (xs: list nat)
  : list nat :=
  match xs with
  | [] => []
  | [x] => [x]
  | _ => let (l,r) :=
           split xs in
         merge (merge_sort l)
               (merge_sort r)
  end.
\end{lstlisting}
\end{block}
\end{columns}
\vfill
\begin{itemize}
  \item Two codebases to write, maintain, and keep in sync
  \item Your \emph{actual} Python is never verified
  \item \jugmark{Any divergence voids the guarantee}
\end{itemize}
\end{frame}

%% Slide 11: Pain Point 3 — Binary Trust
\begin{frame}
\frametitle{Pain Point 3 --- Binary Trust}
\begin{center}
\begin{tikzpicture}[
  box/.style={draw, thick, rounded corners, minimum width=3cm, minimum height=0.8cm,
              font=\small, align=center}
]
  \node[box, fill=jugeoGreen!20]                        (pass) {\textcolor{jugeoGreen}{\textbf{Type-checks}} $\Rightarrow$ fully trusted};
  \node[box, fill=jugeoRed!15, below=0.5cm of pass] (fail) {\textcolor{jugeoRed}{\textbf{Doesn't type-check}} $\Rightarrow$ zero trust};
\end{tikzpicture}
\end{center}
\vspace{0.3cm}
\begin{itemize}
  \item There's no middle ground
  \item What about mixed evidence?
  \begin{itemize}
    \item ``Z3 proved the core math''
    \item ``Tests covered the edge cases''
    \item ``AI suggested the invariant''
    \item ``Manual review checked the IO''
  \end{itemize}
  \item \jugmark{No way to combine or account for these}
\end{itemize}
\end{frame}

%% Slide 12: Pain Point 4 — Opaque Failures
\begin{frame}
\frametitle{Pain Point 4 --- Opaque Failures}
\begin{columns}[T]
\column{0.48\textwidth}
\begin{alertblock}{F* says}
  \texttt{(error) Unknown assertion failed}\\[4pt]
  \footnotesize Z3 timed out. No explanation.\\
  No hint about \emph{what} went wrong.
\end{alertblock}
\column{0.48\textwidth}
\begin{alertblock}{Coq says}
  \texttt{Error: Tactic failure.}\\[4pt]
  \footnotesize Which sub-goal? Why?\\
  Good luck if you're not an expert.
\end{alertblock}
\end{columns}
\vfill
\begin{itemize}
  \item Failures are \textbf{unstructured} --- a string, not data
  \item No machine-readable explanation of \emph{what} failed
  \item No suggestion for \emph{how} to fix it
  \item \jugmark{Debugging proofs is harder than debugging code}
\end{itemize}
\end{frame}

%% Slide 13: Pain Point 5 — Python Effects Not Supported
\begin{frame}
\frametitle{Pain Point 5 --- Python Effects}
\begin{center}
\begin{tikzpicture}[
  effect/.style={draw, thick, rounded corners, fill=jugeoOrange!15,
                 minimum width=2.5cm, minimum height=0.7cm, font=\small}
]
  \node[effect]                          (exc)  {\texttt{try / except}};
  \node[effect, right=0.4cm of exc]      (asy)  {\texttt{async / await}};
  \node[effect, right=0.4cm of asy]      (gen)  {\texttt{yield}};
  \node[effect, below=0.4cm of exc]      (ctx)  {\texttt{with} (context mgr)};
  \node[effect, right=0.4cm of ctx]      (meta) {metaclasses};
  \node[effect, right=0.4cm of meta]     (desc) {descriptors};
\end{tikzpicture}
\end{center}
\vfill
\begin{itemize}
  \item These are \emph{core} Python idioms, not exotic features
  \item LEAN, Coq, F* have no native model for any of them
  \item F* has \emph{some} effect support (monadic), but not Python's flavour
  \item \jugmark{You can't verify what you can't model}
\end{itemize}
\end{frame}

%% Slide 14: Pain Point 6 — AI Is External
\begin{frame}
\frametitle{Pain Point 6 --- AI Is External}
\begin{columns}[T]
\column{0.55\textwidth}
\begin{itemize}
  \item GitHub Copilot for LEAN exists
  \item LLM-based tactic suggestion is growing
  \item But AI is a \textbf{bolt-on}:
  \begin{itemize}
    \item No trust accounting
    \item AI proof step = formal proof step
    \item If the LLM hallucinates a lemma, \\the system accepts or rejects --- nothing in between
  \end{itemize}
\end{itemize}
\column{0.42\textwidth}
\begin{center}
\begin{tikzpicture}[
  box/.style={draw, thick, rounded corners, minimum width=2.4cm,
              minimum height=0.8cm, font=\small, align=center}
]
  \node[box, fill=jugeoLight]                     (ai)  {AI suggests\\proof step};
  \node[box, fill=jugeoLight, below=0.5cm of ai]  (tc)  {Type-checker\\accepts/rejects};
  \node[below=0.3cm of tc, font=\scriptsize, text=jugeoRed, align=center]
    {Same trust as\\human-written proof};
  \draw[-{Stealth}, thick, jugeoBlue] (ai) -- (tc);
\end{tikzpicture}
\end{center}
\end{columns}
\vfill
\begin{itemize}
  \item \jugmark{No way to say ``AI contributed 40\%, Z3 proved 60\%''}
\end{itemize}
\end{frame}

%% Slide 15: Pain Point 7 — Module-Level Only
\begin{frame}
\frametitle{Pain Point 7 --- Module-Level Only}
\begin{center}
\begin{tikzpicture}[
  mod/.style={draw, thick, fill=jugeoLight, minimum width=1.2cm,
              minimum height=0.8cm, font=\scriptsize},
  arr/.style={-{Stealth[length=2mm]}, thick, gray}
]
  \foreach \i in {0,...,4} {
    \foreach \j in {0,...,3} {
      \node[mod] (m\i\j) at (1.6*\i, -1.1*\j) {};
    }
  }
  \draw[arr] (m00) -- (m10);  \draw[arr] (m10) -- (m20);
  \draw[arr] (m01) -- (m11);  \draw[arr] (m21) -- (m31);
  \draw[arr] (m02) -- (m12);  \draw[arr] (m12) -- (m22);
  \draw[arr] (m03) -- (m13);  \draw[arr] (m23) -- (m33);
  \draw[arr] (m30) -- (m40);  \draw[arr] (m31) -- (m41);
  \draw[arr] (m32) -- (m42);  \draw[arr] (m33) -- (m43);
  \draw[arr] (m20) -- (m30);
  \draw[arr] (m22) -- (m32);

  \node[draw=jugeoRed, thick, dashed, fit=(m11), inner sep=3pt, label={[font=\scriptsize,
    text=jugeoGreen]below:verified}] {};

  \node[right=0.4cm of m43, font=\small, text=gray, align=left]
    {Hundreds of files\\One verified?};
\end{tikzpicture}
\end{center}
\vfill
\begin{itemize}
  \item Proof assistants verify \emph{one module at a time}
  \item No built-in mechanism for cross-module reasoning
  \item Real projects have hundreds of interacting files
  \item \jugmark{Verification doesn't scale to whole projects}
\end{itemize}
\end{frame}

%% Slide 16: Pain Point 8 — Proofs Don't Ship With Code
\begin{frame}
\frametitle{Pain Point 8 --- Proofs Don't Ship}
\begin{center}
\begin{tikzpicture}[
  box/.style={draw, thick, rounded corners, minimum width=3cm, minimum height=0.9cm,
              font=\small, align=center},
  arr/.style={-{Stealth[length=3mm]}, thick}
]
  \node[box, fill=jugeoLight]                        (src)  {Coq / F* source\\+ proofs};
  \node[box, fill=jugeoGreen!15, right=2cm of src]   (ext)  {Extracted\\OCaml / C};
  \node[box, fill=jugeoRed!10, right=2cm of ext]     (dep)  {Deployed\\binary};

  \draw[arr, jugeoBlue]  (src) -- node[above, font=\scriptsize]{extraction} (ext);
  \draw[arr, jugeoOrange] (ext) -- node[above, font=\scriptsize]{compile} (dep);

  \node[below=0.2cm of src, font=\scriptsize, text=jugeoGreen] {proofs here};
  \node[below=0.2cm of ext, font=\scriptsize, text=jugeoRed]   {no proofs};
  \node[below=0.2cm of dep, font=\scriptsize, text=jugeoRed]   {no proofs};
\end{tikzpicture}
\end{center}
\vfill
\begin{itemize}
  \item Extracted code carries \textbf{no proof metadata}
  \item Deployed artifact has \textbf{no record} of how it was verified
  \item You must trust the build pipeline, not the artifact
  \item \jugmark{Proofs evaporate at extraction time}
\end{itemize}
\end{frame}

%% Slide 17: Summary of Pain Points
\begin{frame}
\frametitle{Summary --- 8 Pain Points}
\begin{center}
\small
\begin{tabular}{@{} r l l @{}}
  \toprule
  \textbf{\#} & \textbf{Pain Point} & \textbf{Impact} \\
  \midrule
  1 & No Python target              & Largest community excluded \\
  2 & Must rewrite everything       & Double codebase, drift risk \\
  3 & Binary trust (all or nothing) & Can't combine evidence sources \\
  4 & Opaque failures               & Unstructured, unhelpful errors \\
  5 & Python effects unsupported    & Core idioms can't be modeled \\
  6 & AI is a bolt-on               & No trust accounting for LLMs \\
  7 & Module-level only             & Doesn't scale to real projects \\
  8 & Proofs don't ship             & Artifact carries no evidence \\
  \bottomrule
\end{tabular}
\end{center}
\vfill
\begin{center}
\large These aren't minor inconveniences ---\\
they're \jugmark{structural limitations} of the extraction model.
\end{center}
\end{frame}

%% Slide 18: What If There Were a Better Way?
\begin{frame}
\frametitle{What If\ldots?}
\begin{center}
\vspace{0.5cm}
{\Large What if we could\ldots}\\[0.6cm]
\begin{itemize}
  \item[\textcolor{jugeoGreen}{$\checkmark$}] Verify \textbf{Python} directly --- no rewriting
  \item[\textcolor{jugeoGreen}{$\checkmark$}] Write proofs \textbf{in Python} --- no new language
  \item[\textcolor{jugeoGreen}{$\checkmark$}] Mix evidence sources with \textbf{trust accounting}
  \item[\textcolor{jugeoGreen}{$\checkmark$}] Ship proofs \textbf{with} the code --- not separate
  \item[\textcolor{jugeoGreen}{$\checkmark$}] Scale verification to \textbf{entire projects}
\end{itemize}
\vfill
{\large\itshape ``What if proofs traveled with the code?''}
\end{center}
\end{frame}

%% Slide 19: Enter JuGeo — Judgment Geometry
\begin{frame}
\frametitle{Enter JuGeo --- Judgment Geometry}
\begin{center}
\vspace{0.3cm}
{\LARGE \jugmark{JuGeo}}\\[0.5cm]
{\large A \textbf{proof-carrying Python} framework\\[4pt]
where proofs are \textcolor{jugeoGreen}{\textbf{geometric objects}}, not program terms.}
\end{center}
\vfill
\begin{columns}[T]
\column{0.48\textwidth}
\begin{block}{Traditional}
  Proof = a program that type-checks
\end{block}
\column{0.48\textwidth}
\begin{exampleblock}{JuGeo}
  Proof = a geometric object attached to code
\end{exampleblock}
\end{columns}
\vfill
\begin{center}
\small Proofs are \textbf{first-class data}, computed, composed, and shipped alongside your Python.
\end{center}
\end{frame}

%% Slide 20: Roadmap of This Talk
\begin{frame}
\frametitle{Roadmap}
\begin{center}
\begin{tikzpicture}[
  part/.style={draw, thick, rounded corners, minimum width=2.4cm, minimum height=0.68cm,
               font=\scriptsize\bfseries, align=center, text=white},
  done/.style={part, fill=jugeoGreen},
  todo/.style={part, fill=jugeoBlue}
]
  \node[done]                         (p1) {I.\ Motivation\;\checkmark};
  \node[todo, right=0.25cm of p1]     (p2) {II.\ The Big Idea};
  \node[todo, right=0.25cm of p2]     (p3) {III.\ Trust};

  \node[todo, below=0.4cm of p1]      (p4) {IV.\ Verification};
  \node[todo, right=0.25cm of p4]     (p5) {V.\ Proof-Carrying\\Python};
  \node[todo, right=0.25cm of p5]     (p6) {VI.\ Head-to-Head};

  \node[todo, below=0.4cm of p4]      (p7) {VII.\ Architecture};
  \node[todo, right=0.25cm of p7]     (p8) {VIII.\ Ideation};
  \node[todo, right=0.25cm of p8]     (p9) {IX.\ Orchestration};

  \node[todo, below=0.4cm of p8]      (p10) {X.\ Vision};

  \draw[-{Stealth}, thick, jugeoBlue!50] (p1) -- (p2);
  \draw[-{Stealth}, thick, jugeoBlue!50] (p2) -- (p3);
  \draw[-{Stealth}, thick, jugeoBlue!50] (p3) -- (p4);
  \draw[-{Stealth}, thick, jugeoBlue!50] (p4) -- (p5);
  \draw[-{Stealth}, thick, jugeoBlue!50] (p5) -- (p6);
  \draw[-{Stealth}, thick, jugeoBlue!50] (p6) -- (p7);
  \draw[-{Stealth}, thick, jugeoBlue!50] (p7) -- (p8);
  \draw[-{Stealth}, thick, jugeoBlue!50] (p8) -- (p9);
  \draw[-{Stealth}, thick, jugeoBlue!50] (p9) -- (p10);
\end{tikzpicture}
\end{center}
\vfill
\begin{center}
\small Part I complete --- next: \jugmark{The Big Idea} behind Judgment Geometry
\end{center}
\end{frame}
